%!TEX root = ../../thesis.tex

\section{Determinism and the Planning Contract}\label{sec:corpus-determinism}

% TODO: The central design section of this chapter --- give it the space that implies.
%
% The problem: real mutation operators are stochastic, and a stochastic transition
% breaks constraint C4 of \cref{sec:mdp-constraints}, because MCTS plans over a
% learned deterministic latent model.
%
% The resolution: make each operator a deterministic *enumeration* and sweep it over
% the chosen position --- every resulting seed is executed, and the results are
% collapsed into one successor state by a reduce function. State the operators used
% and the sweep width each induces.
%
% Consequences, all worth stating separately:
%   - the transition becomes deterministic, so the planner's model is well posed,
%   - one agent decision now buys up to a sweep-width's worth of target executions,
%     which directly attacks the objection that MCTS overhead cannot be amortized at
%     fuzzing throughput (\cref{sec:rlfuzz-difficulties}) --- quantify it against the
%     throughput numbers of \cref{sec:platform-throughput},
%   - it does NOT multiply training data: the whole sweep still yields one transition
%     for the learner,
%   - therefore the reduce function silently defines what the operator *means* to the
%     agent, and must be chosen deliberately rather than defaulted into.
%
% Frame the whole section as the price of choosing EfficientZero over a model-free
% learner: \cite{bottinger_deep_2018} deliberately *kept* mutation stochastic because
% Q-learning handles stochastic MDPs natively. The same design pressure produces
% opposite decisions. That is the clearest illustration in the thesis of the algorithm
% dictating the fuzzer's design.
